\documentclass[10pt,twoside]{book}
\usepackage[paperwidth=6in,paperheight=9in,inner=0.8in,outer=0.7in,top=0.85in,bottom=0.8in,headsep=0.25in]{geometry}
\usepackage{newtxtext,newtxmath}
\usepackage{amsmath}
\usepackage{microtype}
\usepackage{fancyhdr}
\usepackage{titlesec}
\usepackage{pgfplots}
\pgfplotsset{compat=1.18}
\titleformat{\section}{\normalfont\large\bfseries}{\thesection}{0.8em}{}
\titlespacing*{\section}{0pt}{2.2ex plus .5ex}{1.2ex}
\pagestyle{fancy}
\fancyhf{}
\fancyhead[LE]{\small\thepage\hspace{1.2em}\textsc{Statistical Physics: A First Course}}
\fancyhead[RO]{\small\textsc{3\enspace The Canonical Ensemble}\hspace{1.2em}\thepage}
\renewcommand{\headrulewidth}{0pt}
\setcounter{chapter}{3}
\setcounter{section}{1}
\setcounter{equation}{8}
\renewcommand{\theequation}{3.\arabic{equation}}
\renewcommand{\thesection}{3.\arabic{section}}
\setcounter{page}{41}
\newcommand{\dd}{\mathrm{d}}
\begin{document}

\section{The partition function}

The result of the last section is that a system in equilibrium with a reservoir at temperature $T$ occupies a microstate $s$ of energy $E_s$ with a probability proportional to the Boltzmann factor $e^{-E_s/kT}$. Writing $\beta = 1/kT$ and requiring that the probabilities sum to one gives the \emph{canonical distribution},
\begin{equation}
  P_s = \frac{e^{-\beta E_s}}{Z}, \qquad Z = \sum_s e^{-\beta E_s}.
\end{equation}
The normalizing sum $Z$ is called the \emph{partition function}. The name (from the German \emph{Zustandssumme}, ``sum over states'') undersells it. $Z$ depends on the temperature and, through the energies $E_s$, on the volume and on whatever else fixes the levels of the system, and once it is known as a function of these variables every equilibrium property follows by differentiation. The rest of this chapter is largely an elaboration of that remark.

Two features of the definition deserve comment at once. First, the sum runs over \emph{microstates}, not over energies: if an energy level $E$ is shared by $g(E)$ distinct states, it contributes $g(E)\,e^{-\beta E}$, and in this form the competition between energy and entropy is visible directly. The Boltzmann factor falls with $E$; the degeneracy usually rises steeply with it, and the product is sharply peaked. Second, the zero of energy is arbitrary. Shifting every $E_s$ by a constant $E_0$ multiplies $Z$ by $e^{-\beta E_0}$ and leaves every $P_s$ unchanged, so only energy differences have physical meaning; we shall see that $\ln Z$ shifts by $-\beta E_0$, and the mean energy by $E_0$, exactly as it should.

At low temperature ($\beta \to \infty$) the sum is dominated by the ground state, and $Z \approx g_0 e^{-\beta E_0}$. At high temperature every Boltzmann factor approaches one and $Z$ simply counts the accessible states. Between these limits $Z$ interpolates smoothly, and the way it does so encodes the spectrum of the system.

\section{Average values}

Any quantity $X$ that takes the value $X_s$ in microstate $s$ has the canonical average
\begin{equation}
  \langle X\rangle = \sum_s X_s P_s = \frac{1}{Z}\sum_s X_s\, e^{-\beta E_s}.
\end{equation}
For the energy itself the sum can be done without knowing anything further about the system. Differentiating a single Boltzmann factor with respect to $\beta$ brings down a factor of $-E_s$,
\begin{equation}
  \frac{\partial}{\partial\beta}\, e^{-\beta E_s} = -E_s\, e^{-\beta E_s},
\end{equation}
so the sum $\sum_s E_s e^{-\beta E_s}$ is $-\partial Z/\partial\beta$, and
\begin{equation}
  U \equiv \langle E\rangle = -\frac{1}{Z}\frac{\partial Z}{\partial \beta} = -\frac{\partial \ln Z}{\partial \beta}.
\end{equation}
This is the first of the promised derivatives. Notice that it is $\ln Z$, not $Z$, that appears naturally; we shall meet this again.

The same trick applied twice gives the mean square energy, $\langle E^2\rangle = Z^{-1}\,\partial^2 Z/\partial\beta^2$, and hence the variance of the energy about its mean:
\begin{align}
  \sigma_E^2 = \langle E^2\rangle - \langle E\rangle^2
  &= \frac{1}{Z}\frac{\partial^2 Z}{\partial\beta^2} - \left(\frac{1}{Z}\frac{\partial Z}{\partial\beta}\right)^{\!2} \notag\\
  &= \frac{\partial^2 \ln Z}{\partial \beta^2} = -\frac{\partial U}{\partial \beta}.
\end{align}
Since $\dd\beta = -\dd T/kT^2$, the last derivative is $kT^2\,\partial U/\partial T$, and $\partial U/\partial T$ at fixed volume is the heat capacity. Therefore
\begin{equation}
  \sigma_E^2 = kT^2 C_V .
\end{equation}
The fluctuations of the energy, which the canonical ensemble permits and the microcanonical forbids, are fixed by a quantity we can measure with a calorimeter. Both $U$ and $C_V$ are extensive, proportional to the number of particles $N$, so the relative size of the fluctuations is
\begin{equation}
  \frac{\sigma_E}{U} \propto \frac{\sqrt{N}}{N} = \frac{1}{\sqrt{N}} .
\end{equation}
For a macroscopic system, with $N \sim 10^{23}$, this is about one part in $10^{11}$. The energy of a system in contact with a reservoir is not fixed, but it is fixed in practice, and this is why the canonical and microcanonical ensembles give the same thermodynamics for large systems.

\section{The Helmholtz free energy}

Thermodynamics tells us that the natural potential for a system at fixed $T$, $V$ and $N$ is the Helmholtz free energy $F = U - TS$. Its connection with the partition function is the most compact statement in the subject:
\begin{equation}
  F = -kT \ln Z .
\end{equation}
To verify it, recall that $S = -(\partial F/\partial T)_V$, so that $F = U + T(\partial F/\partial T)_V$. In terms of $\beta$ this reads $U = \partial(\beta F)/\partial\beta$. Comparing with the expression for $U$ above, $\beta F$ and $-\ln Z$ have the same $\beta$-derivative, so they differ at most by a constant independent of temperature; the third law, or a direct check at $T\to 0$ where both reduce to $\beta E_0 - \ln g_0$, fixes the constant at zero.

With $F$ in hand the entropy follows. Substituting $F = -kT\ln Z$ and $U = \sum_s E_s P_s$ into $S = (U-F)/T$, and using $\ln P_s = -\beta E_s - \ln Z$, gives
\begin{equation}
  S = -k \sum_s P_s \ln P_s ,
\end{equation}
the form in which Gibbs wrote it and the one that carries over unchanged to information theory. When all $\Omega$ accessible states are equally likely it reduces to Boltzmann's $S = k\ln\Omega$.

The logarithm explains why $\ln Z$ rather than $Z$ keeps appearing. For two systems that do not interact, every microstate of the pair is a pair of microstates and the energies add, so the partition function of the whole is the product $Z = Z_1 Z_2$. Then $\ln Z$, $F$, $U$ and $S$ are all sums of the contributions of the parts, as extensive quantities must be.

\section{Example: a two-level system}

The simplest system with a nontrivial partition function has just two states, of energies $0$ and $\varepsilon$. It describes, for instance, a spin-$\frac12$ in a magnetic field, or an impurity with one low-lying excited level. Here
\begin{equation}
  Z = 1 + e^{-\beta\varepsilon}, \qquad
  U = -\frac{\partial\ln Z}{\partial\beta} = \frac{\varepsilon}{e^{\beta\varepsilon} + 1}.
\end{equation}
At low temperature the system sits in its ground state and $U \to 0$; at high temperature the two states are equally occupied and $U \to \varepsilon/2$. It never reaches a population inversion: no positive temperature puts more systems in the upper state than in the lower. The heat capacity is
\begin{equation}
  C = \frac{\partial U}{\partial T} = k\,(\beta\varepsilon)^2\,\frac{e^{\beta\varepsilon}}{\left(e^{\beta\varepsilon}+1\right)^2} ,
\end{equation}
which vanishes at both ends and has a single peak near $kT \approx 0.42\,\varepsilon$ (Fig.~3.2). A bump of this shape in a measured heat capacity, the \emph{Schottky anomaly}, is the signature of a gap in a spectrum and a way to measure it.

\begin{center}
\begin{tikzpicture}
\begin{axis}[width=3.6in,height=1.9in,xmin=0,xmax=2.5,ymin=0,ymax=0.5,
  xlabel={$kT/\varepsilon$},ylabel={$C/k$},samples=200,domain=0.02:2.5,
  tick label style={font=\footnotesize},label style={font=\small},axis lines=left]
\addplot[thick] {(1/x)^2*exp(1/x)/(exp(1/x)+1)^2};
\draw[dashed] (axis cs:0.417,0) -- (axis cs:0.417,0.439);
\end{axis}
\end{tikzpicture}

\small\textbf{Figure 3.2}\enspace Heat capacity of a two-level system: the Schottky anomaly.
\end{center}

Of the formulas in this chapter, the one to carry away is the chain of derivatives: from $Z$ to $\ln Z$, to $U$ and $\sigma_E$ by differentiating in $\beta$, and to $F$ and $S$ by multiplying by $-kT$ and differentiating in $T$. Everything in the next three chapters is an application of it.

\section*{Problems}
\small
\noindent\textbf{3.1}\enspace Show that shifting every energy by $E_0$ changes $\ln Z$ by $-\beta E_0$ and $F$ by $E_0$, and leaves $S$ unchanged.

\smallskip
\noindent\textbf{3.2}\enspace For the two-level system, find the entropy $S(T)$ and check its limits as $T\to 0$ and $T\to\infty$.

\smallskip
\noindent\textbf{3.3}\enspace Locate the maximum of the Schottky heat capacity numerically and confirm $kT \approx 0.42\,\varepsilon$.

\smallskip
\noindent\textbf{3.4}\enspace A system has levels $E_n = n\varepsilon$, $n = 0,1,2,\ldots$, each nondegenerate. Sum the partition function and find $U(T)$. Which physical system is this?
\end{document}
